\documentclass[fleqn]{article}
%\textwidth=6in
%\textheight=9in
%\pagestyle{empty}
\usepackage{amsmath,amssymb,amsthm,ragged2e,siunitx,graphicx,multicol,wrapfig,mathtools,parskip}
\usepackage[margin=0.7in]{geometry}
\usepackage[T1]{fontenc}
\usepackage[utf8]{inputenc}
\usepackage[dvipsnames]{xcolor}
\usepackage[export]{adjustbox}
\newtheorem{thm}{{\sc Theorem}}[section]
\newtheorem{prop}[thm]{{\sc Proposition}}
\newtheorem{cor}[thm]{{\sc Corollary}}
\newtheorem{lem}[thm]{{\sc Lemma}}
\newtheorem{df}[thm]{{\sc Definition}}

\DeclarePairedDelimiter\ceil{\lceil}{\rceil}
\DeclarePairedDelimiter\floor{\lfloor}{\rfloor}

\DeclareMathOperator{\sech}{sech}
\DeclareMathOperator{\csch}{csch}

% automatically makes new paragraphs have no indent
\setlength{\parindent}{0pt}

\newcommand{\n}{\noindent}
\newcommand{\ds}{\displaystyle}
\newcommand{\rar}{\rightarrow}

% new command for underscore
\newcommand{\underscore}{\underline{\hspace{1cm}}}

% this is a command for the space between asterisk or plus and the problem statement
\newcommand{\gap}{\hspace{.25cm}}

\newcommand*\colvec[1]{\begin{bmatrix*}[r]#1\end{bmatrix*}} 

% Uncomment this to write vectors bolded 
\renewcommand{\vec}[1]{\boldsymbol{#1}}

% Uncomment this to write vectors bolded AND with arrow on top
\newcommand{\avec}[1]{\overrightarrow{\boldsymbol{#1}}}
% this notation has an arrow on top of vector which is redundant if it is bolded

% Uncomment this to write vectors unbolded with long arrow on top
% \renewcommand{\vec}[1]{\overrightarrow{#1}}

% this is a command for a coefficient matrix with right-aligned entries (good for aligning negative signs)
\newcommand{\mat}[1]{\begin{bmatrix*}[r]#1\end{bmatrix*}} 

% this is a command for a coefficient matrix with center-aligned entries (good for aligning negative signs)
\newcommand{\cmat}[1]{\begin{bmatrix*}[c]#1\end{bmatrix*}} 

% this is a command for a DETERMNINANT matrix with right-aligned entries (good for aligning negative signs) (ABSOLUTE VALUE BARS)
\newcommand{\vmat}[1]{\begin{vmatrix*}[r]#1\end{vmatrix*}} 

% this is a command for black squares (useful for pivots in matrices) when already in math mode (between $ signs)
\newcommand{\bsq}{\blacksquare}






%%  environment for AUGMENTED  matrices
\newenvironment{amatrix}[1]{%
  \left[\begin{array}{@{}*{#1}{r}|r@{}}
}{%
  \end{array}\right]
}


%%% new command for Augmented BAR alone (helpful for putting augmented bar somewhere in middle of matrix)
\newcommand\aug{\fboxsep=-\fboxrule\!\!\!\fbox{\strut}\!\!\!}

% define new color (i.e. pink)
\definecolor{mypink1}{rgb}{0.858, 0.188, 0.478}

% define new color (i.e. medium green)
\definecolor{MediumGreen}{rgb}{0.1, 0.8, 0.65}

% define new color (i.e. dark purple)
\definecolor{DarkPurple}{rgb}{0.6, 0, 0.6}

% this is a command for the size of headers
\newcommand{\header}[1]{\noindent{\Large \textcolor{mypink1}{\textbf{#1}}}}

% this is a command for the size of subheaders
\newcommand{\subheader}[1]{\noindent{\large \textcolor{red}{\textbf{#1}}}}

% this is a command for answers in RED on same line with a space between problem and answer
\newcommand{\sred}[1]{\hspace{1cm} \textcolor{red}{#1}}

% this is a command for answers in RED on same line with NO space between problem and answer
\newcommand{\red}[1]{\textcolor{red}{#1}}

% this is a command for answers in PURPLE on same line with NO space between problem and answer
\newcommand{\DarkPurple}[1]{\textcolor{DarkPurple}{#1}}

% this is a command for answers in GREEN on same line with NO space between problem and answer
\newcommand{\MediumGreen}[1]{\textcolor{MediumGreen}{#1}}

% this is a command for answers in BLUE on same line with a space between problem and answer
\newcommand{\sblue}[1]{\hspace{1cm} \textcolor{blue}{#1}}

% this is a command for answers in BLUE on same line with NO space between problem and answer
\newcommand{\blue}[1]{\textcolor{blue}{#1}}


% this is a command for defining EXAMPLES
\theoremstyle{definition}
\newtheorem{exmp}{Example}[section]




\begin{document}

\centerline{\Large \bf Math Seminar Template Notes}
\centerline{\large \bf Day 10 (Primes + Counting)}
\vspace{\baselineskip}

%\noindent \underline{\hspace{18cm}}

\header{Day 10 (Primes + Counting) } \\


\noindent\fbox{%
\parbox{\textwidth}{
\subheader{Overview} \\ 
%\vspace{0.25cm}
\textbf{\underline{Goals}} 
\begin{enumerate}
\item Prove there are infinitely many primes
\item Solve counting problems
\end{enumerate}
\vspace{.25cm}
}}

\vspace{0.2cm}


\subheader{Proving There are Infinitely Many Primes}

\textbf{Ex 1:} (a) Consider an integer $N$. What is the remainder when you divide $N$ by any of its prime factor, $p_i$? 

\vspace{1.5cm}
(b) What is the remainder when you divide $N+1$ by any of $N$'s prime factors, $p_i$? 

\vspace{3cm}
\textbf{Key Fact:} This means $N+1$ is \underline{\hspace{2cm}} divisible by any of the prime factors of $N$

\vspace{.5cm}

\textbf{Ex 2:} Prove that there are infinitely many prime numbers. 

\vspace{.25cm}
\textbf{Proof Sketch}

\vspace{.25cm}
We will use a proof by contradiction. What will our assumption be? 

\vspace{1.5cm}
If the number of primes is \underline{\hspace{2.5cm}}, we can \underline{\hspace{2cm}} them all: \underline{\hspace{4cm}}

\vspace{1cm}
Let's create a special number using this list that will lead to a contradiction. What number $N$ involves all these primes? (This is very open-ended, so it may be helpful to make multiple guesses)

\vspace{1.5cm}
    \hspace{1cm} Let $N= \underline{\hspace{3cm}}$

\vspace{.25cm}
Using the ideas from the example on the previous page, let's now consider the integer 

\vspace{.25cm}
    \hspace{1cm} \underline{\hspace{1.5cm}}

\vspace{1cm}
{\it \textbf{Q:} If you divide this new number by any prime in our list, what will the remainder be? }

\vspace{1.5cm}
{\it \textbf{Q:} Based on the answer to the previous question, what can we conclude about this new number? How does this lead us to a contradiction? }


%%%
\newpage


\subheader{Counting Problems with Multiplication} 

\textbf{Ex 3:} Three Gryffindors and two Ravenclaws are auditioning for roles as Quidditch commentators. If one Gryffindor and one Ravenclaw must be picked, how many ways are there to select commentators? 

\vspace{.25cm}

\textbf{Soln:} We must make \underline{\hspace{1cm}} \DarkPurple{independent} choices. 

\vspace{.5cm}
	\hspace{1cm} \underline{\hspace{1cm}} \hspace{.25cm} \underline{\hspace{1cm}} $=$ \underline{\hspace{1cm}} ways to choose \underline{\hspace{4.5cm}}

\vspace{3cm}

\textbf{Ex 4*:} (try 2mins*) How many 12-digit passwords can we create if the 1st and 2nd digits must be different, the 3rd digit must be either 0 or 1, and the last two digits must be the same?
% $(x_1 \ x_2 \ x_3) \hspace{.15cm} x_4 \ x_5 \ x_6 \hspace{.1cm} - \hspace{.1cm} x_7 \ x_8 \ x_9 \ x_{10}$ where $x_i \in \{0, 1, 2, \ldots , 9\}$, $x_1, x_4, x_5 \not \in \{0, 1\}$ and $x_2 \in \{0, 1\}$?  

\vspace{.25cm}

\textbf{Soln:} We must make \underline{\hspace{1cm}} \DarkPurple{independent} choices. 

\vspace{4cm}

%%%% something in a box
\noindent\fbox{%
\parbox{\textwidth}{\vspace{.25cm}
\textbf{Key Idea:} \\ 
When we are selecting multiple objects \DarkPurple{independently} (i.e. the number of ways to choose 1st object doesn't affect the 

\vspace{.25cm}
number of ways to choose 2nd object, etc.), then we \underline{\hspace{3cm}} the number of choices for each ``spot"

\vspace{.5cm} 
Suppose we are selecting $k$ objects and there are

\vspace{.25cm}
\begin{itemize}
\item $a_1$ ways to pick the 1st object
\item $a_2$ ways to pick the 2nd object
\item $a_3$ ways to pick the 3rd object
\item \hspace{1cm} $\vdots$
\item $a_k$ ways to pick the $k$th object
\end{itemize}

\vspace{.25cm}
Then, there are \underline{\hspace{3.5cm}} ways to select the $k$ total objects  
\vspace{.25cm}   
}}
%%% end box






%%%
\newpage 



\subheader{Counting Problems with Addition (Cases)} 

\textbf{Ex 5:} A website requires passwords whose characters can be letters, $\{a, b, c, \ldots , z\}$ or digits, $\{0, 1, 2, \ldots , 9\}$. \\ 
The passwords are not case sensitive, so an upper case letter (M) and a lower case letter (m) are treated the same. 

\vspace{.25cm}

(a) How many possible passwords are there that have between 6 to 8 characters, inclusive? 

\vspace{.25cm}

\textbf{Soln:} 

\vspace{5cm}

%%%% something in a box
\noindent\fbox{%
\parbox{\textwidth}{\vspace{.25cm}
\textbf{In general,} when counting the number of elements of $A \cup B \cup C$, if $A, B, C$ are \DarkPurple{pairwise disjoint} (no two of them overlap), then 

\vspace{.25cm}
	\hspace{2cm} $|A \cup B \cup C| = \underline{\hspace{3.5cm}} $

\vspace{.5cm}
Equivalently, this is saying that when we have 

\vspace{.25cm}
completely disjoint \underline{\hspace{2.5cm}}, we must \underline{\hspace{2cm}}


% \vspace{.5cm}
\vspace{.25cm}   
}}
%%% end box

\vspace{.25cm}

\subheader{Counting Problems with Complements} 

\textbf{Ex 5:} A website requires passwords whose characters can be letters, $\{a, b, c, \ldots , z\}$ or digits, $\{0, 1, 2, \ldots , 9\}$. \\ 
The passwords are not case sensitive, so an upper case letter (M) and a lower case letter (m) are treated the same. 

\vspace{.25cm}
(b) How many passwords with 10 characters have at least one ``z" in them?

\vspace{.25cm}

\textbf{Soln:} 

\vspace{4cm}


\textbf{Idea:} Count the \underline{\hspace{3cm}} of what we want and \underline{\hspace{3cm}} from the total number of passwords 

\vspace{.25cm}

$\displaystyle{ \Big| \text{pass. with at least one ``z"} \Big| = }$  


\vspace{2.05cm}
$\displaystyle{ \Big| \text{pass. with at least one ``z"} \Big| = }$ 





%%%
\newpage



\subheader{Principle of Inclusion/Exclusion}

\textbf{Ex 6:} A 20-sided die is rolled twice. In how many ways can the sum of the rolls be 21 or the second roll be even? 

\vspace{.25cm}

It seems like there are \underline{\hspace{3.5cm}} 

\vspace{.5cm}

\textbf{Q:} Is $\displaystyle{ \Big| \text{rolls with sum of 21 or second die is even} \Big| = }$

\vspace{2.5cm}

Let's fix this! 

\vspace{3.5cm}
$\displaystyle{ \Big| \text{rolls with sum of 21 or second die is even} \Big| = }$

\vspace{7cm}


%%%% something in a box
\noindent\fbox{%
\parbox{\textwidth}{\vspace{.25cm}
	\begin{center}
	\textbf{Principle of Inclusion/Exclusion}
	\end{center}

\vspace{.5cm}
If $A, B$ are finite sets, then $|A \cup B| = \underline{\hspace{3.5cm}}$ 
\vspace{3cm}   

In the special case where $A, B$ are \DarkPurple{disjoint}, then $|A \cup B| = \underline{\hspace{2.5cm}}$
\vspace{1cm}
}}
%%% end box

\vspace{.25cm}



%%%
\newpage





\subheader{Practice Problems} 

\red{Basics of Counting}

%%%% something in a box
\noindent\fbox{%
\parbox{\textwidth}{\vspace{.25cm}
	\begin{center}
	\textbf{4 Basic Counting Strategies}
	\end{center}

\vspace{.25cm}
\textbf{(1) Multiplying}

\vspace{.25cm} 
When we are selecting multiple objects \DarkPurple{independently} (i.e. the number of ways to choose 1st object doesn't affect the 

\vspace{.25cm}
number of ways to choose 2nd object, etc.), then we \underline{\hspace{3cm}} the number of choices for each ``spot"

\vspace{.5cm} 
Suppose we are selecting $k$ objects and there are

\vspace{.25cm}
\begin{itemize}
\item $a_1$ ways to pick the 1st object
\item $a_2$ ways to pick the 2nd object
\item $a_3$ ways to pick the 3rd object
\item \hspace{1cm} $\vdots$
\item $a_k$ ways to pick the $k$th object
\end{itemize}

\vspace{.25cm}
Then, there are \underline{\hspace{3.5cm}} ways to select the $k$ total objects  

\vspace{1cm}

\textbf{(2) Adding} 

\vspace{.25cm}
If we are counting the number of ways $A$ could happen \textbf{OR} $B$ could happen, i.e. \underline{\hspace{1.5cm}}, and $A$ and $B$ are

\vspace{.25cm} 
\DarkPurple{disjoint} cases, then $|A \cup B| = \underline{\hspace{2cm}}$. 
\vspace{.25cm}   

\vspace{1.5cm}

\textbf{(3) Complements} 

\vspace{.25cm}
If we are counting the number of stuff with some property, but it is hard to count this directly we could instead use 

\vspace{.25cm} 
the \DarkPurple{complement} which would say 

\vspace{.25cm} 
\hspace{1cm} $| \text{\# of stuff with some property} | =  $

\vspace{3.5cm}

\textbf{(4) Inclusion-Exclusion} 

\vspace{.25cm}
If we are counting the number of ways $A$ could happen \textbf{OR} $B$ could happen, i.e. \underline{\hspace{1.5cm}}, and $A$ and $B$ are

\vspace{.25cm} 
not necessarily disjoint cases, then $|A \cup B| = \underline{\hspace{4cm}}$. 
\vspace{3.5cm}  
}}
%%% end box










%%%
\newpage




\textbf{Ex 1:} Suppose that we need to select a 4-character passcode where each character is a capital English letter. 

\vspace{.25cm}
(a) How many total possible passcodes are possible?

\vspace{3cm} 
% 	\item How many passcodes are there where all 4 characters are different? 

(c) How many passcodes are there if the first and last characters must be the same and the second and third can be anything except they must be different from what the first and last characters are?

\vspace{5cm} 


(d) How many passcodes are there that contain at least one ``A"

\vspace{5.5cm} 
(e)  How many passcodes are there that contain exactly one ``A"


\vspace{4cm} 




%%%
\newpage 





\textbf{Ex 2:} A certain state's license plates consists of 6 characters which can either be capital English letters or digits. How many possible license plates can be made

\vspace{.25cm} 
(a) that consist of three letters followed by three digits or four letters followed by two digits? 

\vspace{5cm}




(b) that begin with a vowel (A, E, I, O, U) or end with a digit?



\vspace{5cm}



\textbf{Ex 3:} In the set up of \#1 on the previous page, how many passcodes are there that 

\vspace{.25cm}
(a) begin or end with the letter ``A"

\vspace{5cm}
(b) contain more than one ``A"









\end{document}








\subheader{What Can We Assume is True and What We Cannot Assume is True?}

\vspace{.25cm}
%%%% something in a box
\noindent\fbox{%
\parbox{\textwidth}{\vspace{.25cm}
\underline{Things we can assume are true} 

\vspace{.25cm}

	\begin{enumerate}
	\item Axioms of real numbers and integers. 

		\begin{itemize}
		\item If $x, y$ are real numbers then $x \pm y$, $xy$, and $\displaystyle{\frac{x}{y}}$ are all real (except if $y=0$ in the division) 

		\vspace{.25cm}

		\item If $x, y$ are integers, then $x \pm y$ and $xy$ are integers. 

		\vspace{.25cm}

		\item \textbf{Distributive Law:} If $x, y, z$ are reals, then $x(y+z) = xy+xz$  

		\vspace{.25cm}

		\item \textbf{Commutative Law:} If $x, y$ are reals, then $x+y=y+x$ and $xy=yx$ 

		\vspace{.25cm}

		\item \textbf{Associative Law:} If $x, y, z$ are reals, then $x+(y+z) = (x+y)+z$. 

		\vspace{.25cm}

		\item \textbf{Transitive Law:} If $x, y, z$ are reals with $x<y$ and $y<z$, then $x<z$. (similarly for other inequalities)
		\end{itemize}

	\vspace{.25cm}
	\item Basic Arithmetic Rules 

		\begin{itemize}
		\item Rules for adding/subtracting/multiplying/dividing fractions (e.g. $\displaystyle{ \frac{a}{b} \cdot \frac{c}{d} = \frac{ac}{bd} }$ ) 

		\vspace{.25cm}

		\item Multiplication by $-1$: e.g. if $a$ is rational, then $-a$ is rational (or if $a$ is an integer, then $-a$ is an integer, etc.) 

		\vspace{.25cm}

		\item Multiplication by $0$: If $a$ is real, then $a \cdot 0 = 0$ 

		\vspace{.25cm}

		\item Let $x, y$ be real. If $xy = 0$, then $x=0$ or $y=0$
		\end{itemize}
	\end{enumerate}
\vspace{.25cm}   
}}
%%% end box


\vspace{.5cm}

%%%% something in a box
\noindent\fbox{%
\parbox{\textwidth}{\vspace{.25cm}
\red{\underline{Examples of things we cannot assume are true}} 

\vspace{.25cm}

	\begin{enumerate}
	\item ``Odd + Odd = Even" \hspace{.5cm} (or other versions of this like ``Odd + Even = Odd") 	

	\vspace{.25cm}

	\item Rational + rational = rational 

	\vspace{.25cm}

	\item Generally any statement more complicated than the ones in the axioms of real numbers and integers 
	\end{enumerate}
\vspace{.25cm}   
}}
%%% end box










%%%% something in a box
\noindent\fbox{%
\parbox{\textwidth}{\vspace{.25cm}
\textbf{Strategies to Show Inductive Step}

\vspace{.25cm}
	\begin{enumerate}
	\item Start with $P(k)$ and do stuff to \underline{\hspace{4cm}} so it becomes $P(k+1)$

	\vspace{.25cm}
	\item Start with one side/part of $P(k+1)$ and use the assumption that \underline{\hspace{4cm}} to get the rest of $P(k+1)$
	\end{enumerate}
\vspace{.25cm}   
}}
%%% end box




\begin{multicols}{5} 
Type 
\columnbreak 

Symbol 
\columnbreak

Definition 
\columnbreak

\hspace{1cm}
 \columnbreak
 
\hspace{1cm}
\end{multicols}






%%%%%%%%%%%%%%%%%%%%%%%%%%%%%%%%%%%%%%%%%%%%%%%%%%%%%%%%%




%image aligned
%%
\includegraphics[scale=0.7, right]{BlankAxes}
\vspace{-10pt}


%image centered
%%
\begin{center}
\includegraphics[scale=1.4]{SinCosTanDrawGraphs}
\end{center} 
\vspace{-20pt}



%%% incorrect/correct show something with 2 approaches
\vspace{.25cm} 
\begin{multicols}{2} 
\red{\underline{Incorrect}} \columnbreak

\MediumGreen{\underline{Correct}} 
\end{multicols}


%%% Scratch work for proof
\hspace{7cm} \MediumGreen{\underline{Scratch Work}}
\begin{multicols}{2} 
\hspace{4cm} \MediumGreen{\underline{Given/Know}} \columnbreak

\MediumGreen{\underline{Want to show}} 
\end{multicols}

\vspace{2.5cm} 



%%%% something in a box
\noindent\fbox{%
\parbox{\textwidth}{\vspace{.25cm}
\textbf{Def:} \\ (1) The \DarkPurple{intersection} of sets $A$ and $B$ is the set of all elements in \textbf{both} $A$ \textbf{and} $B$, denoted $\underline{\hspace{2cm}}$. 
\vspace{.25cm}   
}}
%%% end box


%%% big row reduction arrow
\xrightarrow{\hspace*{2cm}}  



%%% table with extra vertical spacing
\begin{tabular}{ c|c|c } 
$f'(x)$ &  $f(x)$  &  $F(x)$  \\ \hline
\rule{0pt}{3ex}           &  $x^n$  &             \\
 \rule{0pt}{5ex}          &  $\displaystyle{x^{-1}=\frac{1}{x}}$   &             \\
\rule{0pt}{5ex}           &  $e^x$  &             \\
\rule{0pt}{5ex}           &  $e^{kx}$  &             \\
\rule{0pt}{5ex}           &  $g(x) \pm h(x)$  &             \\
\rule{0pt}{5ex}           &  $cg(x)$  &             \\
\rule{0pt}{5ex}           &  $g(x)h(x)$  &             \\
\rule{0pt}{5ex}           &  $\displaystyle{\frac{g(x)}{h(x)}}$  &             \\
\end{tabular}



%%% coefficient matrix template
$A = \mat{ 1 & 0 & 1 & 2 & 2 \\ 0 & 0 & 1 & 0 & 1 \\ 2 & 0 & 1 & 4 & 3 }$


%%% augmented matrix template (the number in brackets indicates which column the augmented bar will be placed AFTER)
$\begin{amatrix}{5}
1 & -1 & 0 & -5 & 0 & -9  \\ 
0 & 0  & 1 &  4 & 0 & 4 \\ 
0 & 0  & 0 &  0 & 1 & 3  
 \end{amatrix}$
 
 
 
 %%%% template for multiple alignment points in system of equations
 $\begin{alignedat}{3}
		x_1 &+ x_2 &&- x_3 &&= -6 \\
		3x_1 &+ x_2 &&+x_3  &&= 8 \\
		4x_1 &+ 2x_2 &&      &&= 2 \\
		\end{alignedat}$
